\documentclass[11pt,a4paper]{article}
\usepackage[utf8]{inputenc}
\usepackage[T1]{fontenc}
\usepackage{amsmath,amssymb,amsthm,mathtools,stmaryrd}
\usepackage{tikz-cd}
\usepackage{enumitem}
\usepackage{listings}
\usepackage{xcolor}
\usepackage{hyperref}
\usepackage[margin=1in]{geometry}
\usepackage{fancyhdr}
\usepackage{longtable}
\usepackage{booktabs}

% ── Theorem environments ──
\theoremstyle{definition}
\newtheorem{definition}{Definition}[section]
\newtheorem{example}[definition]{Example}
\newtheorem{remark}[definition]{Remark}
\newtheorem{notation}[definition]{Notation}
\theoremstyle{plain}
\newtheorem{theorem}[definition]{Theorem}
\newtheorem{lemma}[definition]{Lemma}
\newtheorem{proposition}[definition]{Proposition}
\newtheorem{corollary}[definition]{Corollary}

% ── Listings ──
\lstset{
  language=Python, basicstyle=\ttfamily\small,
  keywordstyle=\color{blue}\bfseries, commentstyle=\color{gray}\itshape,
  stringstyle=\color{red!60!black}, showstringspaces=false,
  breaklines=true, frame=single, numbers=left,
  numberstyle=\tiny\color{gray}, xleftmargin=2em,
}

% ── Macros (matching the CCTT monograph) ──
\newcommand{\den}[1]{\llbracket #1 \rrbracket}
\newcommand{\Path}{\mathsf{Path}}
\newcommand{\Nat}{\mathbb{N}}
\newcommand{\Int}{\mathbb{Z}}
\newcommand{\Bool}{\mathbb{B}}
\newcommand{\Set}{\mathbf{Set}}
\newcommand{\Cat}{\mathbf{Cat}}
\newcommand{\Type}{\mathsf{Type}}
\newcommand{\Fib}{\mathsf{Fib}}
\newcommand{\Cech}{\check{\mathrm{C}}}
\newcommand{\PyObj}{\mathsf{PyObj}}
\newcommand{\fold}{\mathsf{fold}}
\newcommand{\cata}{\mathsf{cata}}
\newcommand{\op}{\mathsf{op}}
\newcommand{\iso}{\cong}
\newcommand{\covers}{\mathcal{U}}
\newcommand{\Sem}{\mathsf{Sem}}
\newcommand{\Iso}{\mathsf{Iso}}
\newcommand{\Hzero}{H^0}
\newcommand{\Hone}{H^1}
\newcommand{\Htwo}{H^2}
\newcommand{\site}{\mathcal{S}}
\newcommand{\presheaf}{\mathcal{F}}
\newcommand{\restr}[2]{#1\!\!\mid_{#2}}
\newcommand{\interval}{\mathbb{I}}
\newcommand{\rec}{\mathsf{rec}}
\newcommand{\elim}{\mathsf{elim}}
\newcommand{\transp}{\mathsf{transp}}
\newcommand{\comp}{\mathsf{comp}}
\newcommand{\hfilling}{\mathsf{hfill}}
\DeclareMathOperator{\Hom}{Hom}
\newcommand{\defeq}{\coloneqq}
\newcommand{\Glue}{\mathsf{Glue}}

\pagestyle{fancy}
\fancyhf{}
\fancyhead[L]{\small\itshape The Synergy Atlas}
\fancyhead[R]{\small\thepage}

% ══════════════════════════════════════════════════════════════════

\title{%
  \textbf{The Complete Synergy Atlas}\\[6pt]
  \Large Every Interaction Between Cubical Type Theory\\
  and Sheaf Cohomology in $\mathrm{C}^4$
}
\author{CCTT Project}
\date{\today}

\begin{document}
\maketitle

\begin{abstract}
Cubical type theory and sheaf cohomology are each powerful
independently.  Cubical type theory gives proof-relevant equality,
transport along paths, and higher-dimensional composition.  Sheaf
cohomology gives local-to-global assembly, decomposition of proof
obligations over open covers, and obstruction theory.

The Cubical Cohomological Calculus of Constructions ($\mathrm{C}^4$)
is not merely the union of these two theories --- it exploits their
\emph{interactions}.  This document is a comprehensive atlas of
\textbf{34 synergies} between the cubical and cohomological
components, organized by structural level (types, proofs, covers,
search, metatheory, Python semantics, trust).  Each synergy is
classified by whether it provides new \textsc{Expressivity},
\textsc{Tractability}, or \textsc{Automation}.

We prove that the verification power of $\mathrm{C}^4$ is strictly
greater than the union of its components.
\end{abstract}

\tableofcontents
\newpage

% ── Classification key ──
\noindent
We classify each synergy along three axes:

\begin{center}
\begin{tabular}{lp{10cm}}
\toprule
\textbf{Tag} & \textbf{Meaning} \\
\midrule
\textsc{Expressivity} &
  Can state or prove things that neither cubical TT nor sheaf
  cohomology alone can. \\
\textsc{Tractability} &
  Makes an existing proof obligation computationally cheaper
  (fewer VCs, smaller formulas, lower Z3 complexity). \\
\textsc{Automation} &
  Reduces human annotation burden by synthesizing proof
  components or search strategies automatically. \\
\bottomrule
\end{tabular}
\end{center}


% ═══════════════════════════════════════════════════════════════════
\section{Level 0: The Foundational Identification}
\label{sec:syn-foundation}

\subsection{Synergy 0: Descent = HComp}
\label{sec:syn-descent-hcomp}

\textsc{Tags:} Expressivity, Tractability.

\medskip\noindent
\textbf{Cubical side.}
Homogeneous composition ($\hfilling$) fills an $n$-dimensional cube
from compatible face data.  Given partial proofs on the boundary of
$\interval^n$, the Kan condition guarantees a filler:
\[
  \frac{
    \forall k.\; p_k : \Path_{A}\bigl(\partial^0_k(\vec{t}),\;
    \partial^1_k(\vec{t})\bigr)
    \qquad
    \text{edges agree}
  }{
    \hfilling^{i_1,\ldots,i_n}\,A\,[\vec{p}]\,t_{\mathrm{base}}
    : A
  }
\]

\noindent
\textbf{Cohomological side.}
$\Cech$ descent assembles a global section from local sections on
an open cover $\covers = \{U_i\}$ with overlap compatibility:
\[
  \frac{
    s_i \in \presheaf(U_i) \quad
    \restr{s_i}{U_i \cap U_j} = \restr{s_j}{U_i \cap U_j}
  }{
    s \in \presheaf(X) \text{ with } \restr{s}{U_i} = s_i
  }
\]

\noindent
\textbf{The synergy.}
These are the \emph{same} universal construction:
\begin{enumerate}
\item Each fiber $\varphi_k$ of the refinement cover corresponds to
  a cubical face $[i_k = \varepsilon]$.
\item The local proof $p_k$ on fiber $\varphi_k$ is the partial
  boundary datum on that face.
\item The overlap compatibility $\restr{p_i}{\varphi_i \wedge \varphi_j}
  = \restr{p_j}{\varphi_i \wedge \varphi_j}$ is the cubical edge
  agreement condition.
\item The Kan filling condition (``$\hfilling$ exists'') is
  $\Hone(\covers, \presheaf) = 0$ for the $\Cech$ nerve.
\item For propositional coefficients (path types over decidable
  equality --- our setting by Hedberg's theorem), $\Hone = 0$
  automatically, so the filling always exists.
\end{enumerate}

\begin{theorem}[Operational unification]
A single verification-condition generator
\texttt{\_compile\_cubical\_filling} handles both
\textsf{RefinementDescent} and \textsf{HComp} proof terms,
differing only in whether faces carry Z3-decidable refinement
predicates (descent) or pure cubical names (hcomp).
\end{theorem}

\noindent
\textbf{Why neither alone suffices.}
\begin{itemize}
\item Pure cubical TT has no notion of ``branching code structure.''
  It can fill cubes, but it does not know that the cube's faces
  correspond to program branches.
\item Pure sheaf theory has no notion of ``proof dimension.''  It
  can decompose by opens, but has no path structure connecting
  the local proofs.
\item The identification lets the compiler \emph{decompose by branches}
  (cohomological) and \emph{compose proof paths} (cubical) in a
  single pass.
\end{itemize}


\subsection{Synergy 0b: The Two Dimension Systems}
\label{sec:syn-two-dimensions}

\textsc{Tags:} Expressivity.

\medskip\noindent
$\mathrm{C}^4$ has \emph{two independent dimension systems}:

\begin{center}
\begin{tabular}{lll}
\toprule
& \textbf{Cubical dimensions} & \textbf{Cohomological dimensions} \\
\midrule
Parameterized by & Interval variables $i, j, k : \interval$ &
  Refinement predicates $\varphi(x)$ \\
0-cells & Programs & Input values \\
1-cells & Equality proofs (paths) & Fiber restrictions \\
2-cells & Proof coherence (squares) & Cover overlaps \\
Composition & $\hfilling$ (Kan) & Descent (gluing) \\
Transport & Along type paths & Along predicate implications \\
Obstruction & Kan condition & $\Hone = 0$ \\
Decidability & Structural & Z3 \\
\bottomrule
\end{tabular}
\end{center}

A proof over a function with $n$ branches and $m$ proof steps
lives in $\interval^m \times \site^n$ --- an $(m{+}n)$-dimensional
space.  The cubical dimensions parameterize \emph{how} we prove;
the cohomological dimensions parameterize \emph{what} (which inputs)
we prove for.  The full power of $\mathrm{C}^4$ comes from
working in this product space.


% ═══════════════════════════════════════════════════════════════════
\section{Level 1: Type-Level Synergies}
\label{sec:syn-types}


\subsection{Synergy 1: Refinement Types as Cubical Faces}
\label{sec:syn-ref-face}

\textsc{Tags:} Expressivity, Tractability.

\medskip\noindent
A refinement type $\{x : T \mid \varphi(x)\}$ defines both:
\begin{itemize}
\item A \textbf{subtype} in the type-theoretic sense (elements
  satisfying $\varphi$).
\item A \textbf{cubical face} in the proof cube (restricting the
  proof to inputs satisfying $\varphi$).
\end{itemize}

This dual role means that \emph{every type-theoretic refinement
is simultaneously a proof decomposition point}.  Adding a type
annotation to a function parameter automatically provides a
cover for decomposition.

\begin{example}
Annotating \texttt{def f(x: Positive[int])} gives:
\begin{enumerate}
\item A type-level guarantee that $x > 0$ within the body.
\item A cubical face $[\varphi = (x > 0)]$ for the proof.
\item A sheaf restriction $\restr{\presheaf}{\{x > 0\}}$ for local
  verification.
\end{enumerate}
All three are the same data, viewed from different angles.
\end{example}


\subsection{Synergy 2: Univalence in the Refinement Lattice}
\label{sec:syn-ua-ref}

\textsc{Tags:} Expressivity.

\medskip\noindent
Cubical univalence says: if $A \simeq B$ (as types), then
$\Path_\Type(A, B)$.  Applied to the refinement lattice:

\[
  \frac{
    \varphi \Leftrightarrow \psi
    \quad\text{(Z3 proves equivalence)}
  }{
    \Path_\Type\bigl(\{x \mid \varphi\},\; \{x \mid \psi\}\bigr)
  }
\]

This means: \emph{logically equivalent refinement predicates define
the same type}, up to a path.  Proofs on $\{x \mid \varphi\}$ can be
freely used for $\{x \mid \psi\}$ via transport.

\textbf{Why this needs both theories.}
\begin{itemize}
\item The equivalence $\varphi \Leftrightarrow \psi$ is a
  cohomological fact (Z3-decidable predicate logic on the site).
\item The path $\Path_\Type(A, B)$ is a cubical fact (univalence
  axiom).
\item Together, they give \textbf{Z3-decidable type equalities}:
  the compiler automatically identifies refinement types with
  logically equivalent predicates.
\end{itemize}


\subsection{Synergy 3: Fibered Universe Structure}
\label{sec:syn-fibered-universe}

\textsc{Tags:} Expressivity.

\medskip\noindent
The universe $\Type$ in $\mathrm{C}^4$ is \emph{fibered} over the
refinement site:
\[
  \Type \;\to\; \site
\]
Each type $A$ sits over a refinement predicate $\varphi$ that
describes the inputs for which $A$ is defined.  The fiber of
$\Type$ over $\varphi$ is the category of types restricted to
$\{x \mid \varphi\}$.

\textbf{Cubical implication:} paths in $\Type$ are \emph{fibered
paths} --- they must respect the projection to $\site$.  A path
between types $A$ and $B$ exists only if their refinement predicates
are related (by implication or equivalence).

\textbf{Cohomological implication:} type families over the site
are \emph{sheaves of types}.  The gluing condition means that
locally defined types that agree on overlaps can be assembled into
a global type --- this is the \textsf{Glue} type of cubical TT,
now fibered over the refinement cover.

\begin{remark}
This fibered structure is what prevents ``type confusion'' in
dynamically-typed Python: each branch of an \texttt{isinstance}
dispatch lives in a \emph{different fiber} of the universe, and
proofs cannot accidentally leak between fibers without an explicit
transport.
\end{remark}


\subsection{Synergy 4: Sheafification as Type Completion}
\label{sec:syn-sheafification}

\textsc{Tags:} Expressivity, Automation.

\medskip\noindent
Given a presheaf $\presheaf$ on the site (a type family that does
not satisfy the gluing condition), sheafification produces the
``best'' sheaf $\presheaf^+$ that does.  In $\mathrm{C}^4$ terms:

\textbf{Cohomological reading:} sheafification completes a partial
type assignment (e.g., a function defined on some but not all
branches) to a total one.

\textbf{Cubical reading:} sheafification fills in missing faces
of a partial cube --- exactly the $\hfilling$ operation applied to
the type level.

\textbf{Practical consequence:} if a function handles some
\texttt{isinstance} cases but not others, the compiler can
\emph{sheafify} the type: automatically extend the type to cover
missing cases by filling the cube with $\bot$ (runtime error)
or by synthesizing default behavior.


% ═══════════════════════════════════════════════════════════════════
\section{Level 2: Proof-Level Synergies}
\label{sec:syn-proofs}


\subsection{Synergy 5: Transport Along Refinement Paths}
\label{sec:syn-transport-ref}

\textsc{Tags:} Tractability, Automation.

\medskip\noindent
Cubical transport moves proofs along type paths:
$\transp^i\,A\,\varphi\,t : A(1)$ given $t : A(0)$.

Applied to refinement types: if we have a proof $p$ that a function
$f$ satisfies a property under hypothesis $\varphi_1$, and a path
(implication) $\varphi_1 \Leftarrow \varphi_2$, then
$\transp\,p$ gives a proof of the same property under the
\emph{stronger} hypothesis $\varphi_2$.

\[
  \frac{
    p : P\;\text{under}\;\varphi_1
    \qquad
    \varphi_2 \Rightarrow \varphi_1 \;\text{(Z3)}
  }{
    \transp(p) : P\;\text{under}\;\varphi_2
  }
\]

\textbf{Why this needs both.}
\begin{itemize}
\item The direction of transport (from weaker to stronger hypothesis,
  or vice versa depending on variance) is a cubical operation.
\item The path $\varphi_2 \Rightarrow \varphi_1$ is a cohomological
  fact (morphism in the site category), checked by Z3.
\item The synergy: transport \emph{along morphisms in the site}
  lets us reuse proofs across different refinement conditions
  without re-proving.
\end{itemize}

\begin{example}[Transport cascade]
Given:
\begin{itemize}
\item $p : $ ``\texttt{abs(x) >= 0}'' under $\{x > 0\}$
\item Z3 proves: $(x > 0) \Rightarrow (x \geq 0)$ and
  $(x \geq 0) \Rightarrow (x \geq -1)$
\end{itemize}
Two transports give ``\texttt{abs(x) >= 0}'' under $\{x \geq 0\}$
and $\{x \geq -1\}$ --- for free.
\end{example}


\subsection{Synergy 6: Restricted Hypothesis Pushing}
\label{sec:syn-hypothesis}

\textsc{Tags:} Tractability.

\medskip\noindent
When compiling a fiber proof under predicate $\varphi_k$, the
predicate itself should be available as a \emph{hypothesis} in the
Z3 context.  This is:

\begin{itemize}
\item \textbf{Cubically:} evaluating the proof at face $[i_k = 0]$,
  which restricts all terms to that face.
\item \textbf{Cohomologically:} restricting the sheaf to the open
  set $U_k = \{x \mid \varphi_k(x)\}$, under which $\varphi_k$
  is tautologically true.
\end{itemize}

\textbf{Practical impact:} without hypothesis pushing, a fiber proof
of ``$x \geq 0$ under $\{x > 0\}$'' would require an explicit
proof of $x > 0 \Rightarrow x \geq 0$.  With hypothesis pushing,
the assumption $x > 0$ is in the Z3 context, so the implication
is automatic.

This turns descent from mere \emph{bookkeeping} (split by cases,
prove each case from scratch) into genuine \emph{simplification}
(split by cases, prove each case with the case hypothesis freely
available).


\subsection{Synergy 7: Overlap Compatibility = Cubical Edge Agreement}
\label{sec:syn-overlap-edge}

\textsc{Tags:} Expressivity.

\medskip\noindent
The $\Cech$ overlap condition says: on $U_i \cap U_j$, the local
sections must agree.  The cubical edge condition says: on the shared
edge of two faces, the partial proofs must agree.

In $\mathrm{C}^4$, an ``overlap'' is $\varphi_i \wedge \varphi_j$.
Whether this is non-empty is decided by Z3 (cohomological).  If it
is non-empty, the agreement condition is a path equality
(cubical):
\[
  \restr{p_i}{\varphi_i \wedge \varphi_j}
  \;=_{\Path}\;
  \restr{p_j}{\varphi_i \wedge \varphi_j}
\]

\textbf{Neither theory alone provides this.}
\begin{itemize}
\item Pure cubical TT checks face agreement but cannot determine
  which faces overlap (it does not have the Z3-decidable
  predicate structure).
\item Pure sheaf theory requires agreement on overlaps but has no
  notion of ``proof path'' to state the agreement as an equality.
\end{itemize}


\subsection{Synergy 8: Automatic Overlap Synthesis}
\label{sec:syn-auto-overlap}

\textsc{Tags:} Automation.

\medskip\noindent
On non-empty overlaps $\varphi_i \wedge \varphi_j$, compatibility
proofs can often be \emph{synthesized automatically} by combining
cubical and cohomological reasoning:

\begin{enumerate}
\item \textbf{Disjointness (Z3):} if $\varphi_i \wedge \varphi_j
  \equiv \bot$, no compatibility proof is needed.  This is the most
  common case for \texttt{if/elif/else} chains.
\item \textbf{Containment (Z3 + transport):} if $\varphi_i \Rightarrow
  \varphi_j$, then the $j$-proof restricted to the overlap equals
  the $i$-proof transported to the overlap.  Compatibility is
  automatic.
\item \textbf{Refl $\cap$ Refl (kernel):} if both fiber proofs
  are reflexivity, their restrictions to the overlap are both
  reflexivity.  Trivially compatible.
\item \textbf{Z3 discharge:} if the compatibility obligation is a
  quantifier-free formula, discharge it directly to Z3 without
  requiring the user to provide an explicit proof.
\end{enumerate}

\textbf{Synergy mechanism:} strategies 1--2 use the cohomological
site structure (Z3 on predicates); strategy 3 uses cubical
reflexivity; strategy 4 uses both.  The interaction means that
\emph{most overlaps in typical Python code are handled without
manual proofs}.


\subsection{Synergy 9: GluePath --- Local Cubical Paths via Descent}
\label{sec:syn-gluepath}

\textsc{Tags:} Expressivity.

\medskip\noindent
The \textsf{GluePath} proof term glues local \emph{paths}
(not just equalities) via $\Cech$ descent.  Each local proof
varies over the cubical interval \emph{and} is restricted to a
refinement fiber:
\[
  p_k : \interval \;\to\;
  \Path\bigl(\restr{f}{\varphi_k},\;\restr{g}{\varphi_k}\bigr)
\]

The descent condition glues these into a global path
$p : \Path(f, g)$ that works on all inputs.

\textbf{This is strictly more powerful than either component:}
\begin{itemize}
\item Cubical TT alone can build paths, but cannot decompose
  them by input cases.
\item Sheaf cohomology alone can decompose by cases, but cannot
  parameterize the local data over a proof interval.
\item \textsf{GluePath} parameterizes the $\Cech$ data over the
  cubical interval, giving a ``moving'' descent.
\end{itemize}

In the spectral sequence: \textsf{GluePath} computes the $E_1$
page, gluing local $\Hzero$ path-sections into global $\Hzero$.


% ═══════════════════════════════════════════════════════════════════
\section{Level 3: Cover-Level Synergies}
\label{sec:syn-covers}


\subsection{Synergy 10: Cover Algebra as Cubical Face Algebra}
\label{sec:syn-cover-face}

\textsc{Tags:} Tractability.

\medskip\noindent
The refinement site has a rich algebra of covers:
\begin{align*}
  \text{Refinement:}\quad & \covers' \leq \covers
    \;\;\text{iff each } \varphi'_j \text{ refines some } \varphi_i \\
  \text{Product:}\quad & \covers \times \mathcal{V}
    = \{\varphi_i \wedge \psi_j\} \\
  \text{Join:}\quad & \covers \vee \mathcal{V}
    = \{\varphi_i\} \cup \{\psi_j\} \\
  \text{Restriction:}\quad & \restr{\covers}{\psi}
    = \{\varphi_i \wedge \psi\}
\end{align*}

Each of these has a cubical interpretation:
\begin{itemize}
\item \textbf{Refinement} = subdividing faces of a cube into
  smaller faces ($\sim$ cubical subdivision).
\item \textbf{Product} = taking the product of two cubes
  ($\interval^m \times \interval^n = \interval^{m+n}$).
\item \textbf{Join} = taking the union of face systems
  ($\sim$ adding dimensions).
\item \textbf{Restriction} = evaluating a face map
  ($\partial^\varepsilon_k : \interval^n \to \interval^{n-1}$).
\end{itemize}

All four operations are Z3-decidable (cohomological) and correspond
to cubical constructions.  No other type theory has this dual
algebraic structure on proof decompositions.


\subsection{Synergy 11: Product Covers for Interprocedural Composition (K\"unneth)}
\label{sec:syn-kunneth}

\textsc{Tags:} Expressivity, Tractability.

\medskip\noindent
When function $f$ (proved over cover $\covers = \{\varphi_i\}$)
calls function $g$ (proved over cover $\mathcal{V} = \{\psi_j\}$),
the interprocedural proof uses the product cover:
\[
  \covers \times \mathcal{V} = \{\varphi_i \wedge \psi_j\}
\]

The K\"unneth formula for $\Cech$ cohomology gives:
\[
  \Hzero(\covers \times \mathcal{V}, \presheaf)
  \;\cong\;
  \Hzero(\covers, \presheaf) \otimes \Hzero(\mathcal{V}, \presheaf)
\]

\textbf{Practical meaning:} the composite proof decomposes into
\emph{pairs} of fiber proofs.  If $f$ has $m$ branches and $g$
has $n$ branches, the monolithic approach needs $O(2^{m+n})$ in the
worst case; the product cover needs $O(m \cdot n)$ VCs, and if the
covers are independent (many branches don't interact), many product
fibers are empty (Z3 detects this), reducing further.

\textbf{Cubical reading:} the product cover corresponds to the
product cube $\interval^m \times \interval^n$.  The composite
proof fills this product cube from the individual cube-fillings
of $f$ and $g$.


\subsection{Synergy 12: Cover Morphisms and Kan Extension}
\label{sec:syn-kan}

\textsc{Tags:} Tractability, Automation.

\medskip\noindent
A \textbf{cover morphism} $\mu : \mathcal{V} \to \covers$ is a map
where each $\mathcal{V}$-fiber refines some $\covers$-fiber:
$\forall j.\;\exists i.\;\psi_j \Rightarrow \varphi_i$.

Given a proof over the coarse cover $\covers$, the \textbf{left Kan
extension} along $\mu$ produces a proof over the fine cover
$\mathcal{V}$:
\[
  \mathrm{Lan}_\mu : \presheaf(\covers) \to \presheaf(\mathcal{V})
\]

\textbf{Cubically:} this is transport along the morphism path in the
site category.  The proof on each $\mathcal{V}$-fiber is obtained by
restricting the proof on the corresponding $\covers$-fiber.

\textbf{Cohomologically:} Kan extension is the left adjoint of the
restriction functor $\mu^* : \presheaf(\covers) \to \presheaf(\mathcal{V})$.

\textbf{Practical consequences:}
\begin{enumerate}
\item \textbf{Proof reuse under refactoring:} if you proved a
  function correct with a 3-branch cover, and refactoring splits
  one branch into two sub-branches, the Kan extension gives you a
  proof over the 4-branch cover for free (after verifying the cover
  morphism, which is a Z3 check).
\item \textbf{Coarsening for efficiency:} if a fine-grained proof
  is expensive, prove over a coarse cover and let Kan extension
  specialize to the finer cover.
\end{enumerate}


\subsection{Synergy 13: Connection-Based Cover Simplification}
\label{sec:syn-connections}

\textsc{Tags:} Tractability.

\medskip\noindent
The cubical interval $\interval$ has \textbf{connections}
$\wedge_i$ (min) and $\vee_i$ (max), satisfying the De Morgan
algebra laws:
\[
  \neg(i \wedge j) = \neg i \vee \neg j, \qquad
  \neg(i \vee j) = \neg i \wedge \neg j
\]

The refinement lattice has its own $\wedge$ (conjunction) and
$\vee$ (disjunction).  The \textbf{synergy}: these two lattice
structures interact to simplify covers via standard lattice
identities:

\begin{enumerate}
\item \textbf{Absorption:} $\varphi \wedge (\varphi \vee \psi) = \varphi$
  --- collapse nested covers.
\item \textbf{Distribution:}
  $\varphi \wedge (\psi \vee \chi) = (\varphi \wedge \psi)
  \vee (\varphi \wedge \chi)$ --- factor product covers.
\item \textbf{Idempotence:} $\varphi \wedge \varphi = \varphi$
  --- deduplicate fibers.
\item \textbf{Complementation:} $\varphi \wedge \neg\varphi = \bot$
  --- detect dead branches (Z3).
\end{enumerate}

These simplifications reduce the number of VCs generated before
any Z3 solving occurs.  The cubical De Morgan laws guarantee that
the simplifications are \emph{homotopy-invariant} (they preserve
the homotopy type of the nerve).


\subsection{Synergy 14: Nerve Topology as Proof Connectivity}
\label{sec:syn-nerve}

\textsc{Tags:} Expressivity.

\medskip\noindent
The $\Cech$ nerve of a cover $\covers$ is the simplicial complex
whose $k$-simplices are $(k{+}1)$-fold overlaps.  Its topology
encodes the ``shape'' of the proof decomposition.

\textbf{Cohomological reading:} the nerve's $\Hone$ measures
obstructions to gluing.  For propositional coefficients,
$\Hone = 0$ always (Hedberg), but for \emph{non-propositional}
coefficient sheaves (e.g., proof-relevant transport data),
$\Hone$ can be non-trivial.

\textbf{Cubical reading:} the nerve embeds into a cubical set
(via the standard nerve--cubical adjunction).  Its homotopy type
determines which $\hfilling$ operations are possible.

\textbf{Synergy:} a disconnected nerve (components that don't
overlap) means the proof \emph{factorizes} into independent
sub-proofs.  Each connected component of the nerve can be filled
independently.  This is \textbf{proof parallelism} --- detected
cohomologically, executed cubically.


% ═══════════════════════════════════════════════════════════════════
\section{Level 4: Proof Search Synergies}
\label{sec:syn-search}


\subsection{Synergy 15: Spectral Sequence as Proof Search Strategy}
\label{sec:syn-spectral}

\textsc{Tags:} Automation.

\medskip\noindent
The $\Cech$-to-derived spectral sequence:
\[
  E_1^{p,q} = \Cech^p(\covers, \mathcal{H}^q(\presheaf))
  \;\Rightarrow\;
  H^{p+q}(X, \presheaf)
\]
gives a systematic proof search strategy:

\begin{center}
\begin{tabular}{llp{6.5cm}}
\toprule
\textbf{Page} & \textbf{Cubical} & \textbf{Cohomological} \\
\midrule
$E_0$ & --- &
  Extract branch structure from code AST; construct
  refinement cover $\covers$ \\
$E_1$ & Compile each fiber proof (path construction) &
  Local proofs on each open $U_i$; local $\Hzero$ \\
$E_2$ & Check edge agreement &
  Overlap compatibility; differential $d_1$ \\
$E_\infty$ & Kan filling (assemble global path) &
  Gluing into global section \\
\bottomrule
\end{tabular}
\end{center}

\textbf{The key synergy:} the spectral sequence page structure
gives a \emph{priority ordering} for proof search.  If $E_1$ fails
(a fiber proof fails), there is no point checking $E_2$ (overlaps).
If $E_1$ succeeds but $E_2$ fails, the overlap counterexample
identifies exactly which interaction between branches is problematic.

This is strictly more structured than either:
\begin{itemize}
\item Cubical proof search (try compositions until one works ---
  no systematic decomposition).
\item Cohomological resolution (compute derived functors ---
  overkill for decidable coefficient sheaves).
\end{itemize}


\subsection{Synergy 16: CEGAR via Cover Refinement}
\label{sec:syn-cegar}

\textsc{Tags:} Automation.

\medskip\noindent
\textbf{C}ounter\textbf{E}xample-\textbf{G}uided
\textbf{A}bstraction \textbf{R}efinement, adapted to the
cohomological setting:

\begin{enumerate}
\item Start with a coarse cover (e.g., $\{x : \mathrm{int}\}$
  --- single fiber, no decomposition).
\item Try to prove the property on each fiber.
\item If a fiber proof fails, Z3 provides a \textbf{counterexample}
  (a specific input $x = c$ that violates the property under that
  fiber's hypothesis).
\item \textbf{Split the fiber} at the counterexample:
  $\varphi_k \mapsto (\varphi_k \wedge x \leq c),\;
  (\varphi_k \wedge x > c)$.
\item Retry with the refined cover.
\item Iterate until all fibers pass or the maximum depth is reached.
\end{enumerate}

\textbf{Cubical reading:} the counterexample identifies a
\emph{point on a cubical face} where the proof fails.  The
refinement adds a new cubical dimension, splitting the face into
two sub-faces that separate the failing case.

\textbf{Cohomological reading:} this is refinement of the cover
in the site topology.  The refined cover is a finer open covering
whose $\Cech$ cohomology may be simpler (smaller local obligations).

\textbf{Synergy:} the counterexample (Z3, cohomological data)
drives the cubical decomposition.  Neither theory alone provides
this loop: cubical TT has no cover structure to refine, and sheaf
theory has no solver to produce counterexamples.


\subsection{Synergy 17: Obstruction-Guided Proof Repair}
\label{sec:syn-obstruction}

\textsc{Tags:} Automation.

\medskip\noindent
When the spectral sequence does not degenerate at $E_2$ (i.e.,
there are non-trivial differentials), the obstructing cocycle
$d_1(s) \in \Cech^1(\covers, \presheaf)$ tells you
\emph{exactly which overlaps fail}.

\textbf{Cubically:} a non-trivial $d_1$ means two cubical faces
disagree on a shared edge.  The disagreement is a specific pair
of face values that do not match.

\textbf{Cohomologically:} the cocycle class $[d_1(s)]$ in $\Hone$
is the obstruction.  If it is a coboundary ($d_1(s) = \delta^0 t$),
then adjusting the local sections by $t$ makes them compatible.

\textbf{Practical use:} when an overlap compatibility check fails,
the compiler can:
\begin{enumerate}
\item Identify the specific overlap $(U_i \cap U_j)$ that fails.
\item Examine whether the failure is a coboundary (can be fixed by
  adjusting the local proofs).
\item If so, automatically \emph{repair} the local proofs to achieve
  compatibility --- a form of proof synthesis guided by cohomological
  obstruction theory.
\end{enumerate}


% ═══════════════════════════════════════════════════════════════════
\section{Level 5: Compositional Synergies}
\label{sec:syn-compositional}


\subsection{Synergy 18: Transport-Descent Functoriality}
\label{sec:syn-transp-descent}

\textsc{Tags:} Tractability, Automation.

\medskip\noindent
Transport and descent commute:
\[
  \transp\bigl(\mathrm{desc}(\covers, \{p_i\})\bigr)
  \;=\;
  \mathrm{desc}\bigl(\transp(\covers), \{\transp(p_i)\}\bigr)
\]

\textbf{What this means:} if you have a descent proof of
$f$ being correct, and a path $q : f \equiv g$ (e.g., from a
refactoring or library equivalence), then $g$ is correct via
descent over the \emph{same cover structure}, with each fiber
proof transported.

\textbf{Why both theories are needed:}
\begin{itemize}
\item Transport is cubical (moving along a path in the type/program
  space).
\item Descent is cohomological (decomposition over a cover in the
  site).
\item The functoriality says the two structures are \emph{compatible}:
  transport preserves the decomposition.
\end{itemize}

\textbf{Practical consequence:} when a library function is updated
(new version with different implementation but same specification),
the existing proof transfers automatically.  The cost is:
1 path proof (old $\equiv$ new) + 0 re-verification of fiber proofs.


\subsection{Synergy 19: Descent Naturality (Proof Independence of Cover)}
\label{sec:syn-naturality}

\textsc{Tags:} Tractability.

\medskip\noindent
For any cover morphism $\mu : \mathcal{V} \to \covers$, descent
is \emph{natural}:
\[
  \mu^*\bigl(\mathrm{desc}_\covers(s)\bigr)
  \;=\;
  \mathrm{desc}_\mathcal{V}\bigl(\mu^*(s)\bigr)
\]

This means the result of descent does not depend on the choice of
cover (up to the Kan condition).  In cubical terms: the filled cube
is unique (up to homotopy) regardless of how the faces were
decomposed.

\textbf{Practical impact:} you can choose the most \emph{convenient}
cover for proving (the one that makes the Z3 obligations simplest),
knowing that the result is cover-independent.  Different
cover choices yield equivalent global proofs.


\subsection{Synergy 20: Fiber-Preserving Homotopies (Relative Paths)}
\label{sec:syn-relative-paths}

\textsc{Tags:} Expressivity.

\medskip\noindent
A \textbf{fiber-preserving homotopy} is a path in the
proof space that respects the fibration over the refinement site:
\[
  H : \interval \to \prod_{i} \Path\bigl(
  \restr{f}{\varphi_i},\;\restr{g}{\varphi_i}\bigr)
\]
such that the overlap conditions hold for all $t : \interval$.

This is a ``moving descent'' --- a descent proof that varies
continuously over the cubical interval.  It corresponds to the
$E_1$ page of the spectral sequence: a family of $\Cech$ $0$-cochains
parameterized by the interval, satisfying the cocycle condition
at each parameter value.

\textbf{Use case:} proving that two algorithms are ``homotopically
the same'' --- not just equal on outputs, but equal in a way
that is \emph{uniform} over all input branches.  This rules out
proofs that work ``by accident'' on some branches.


\subsection{Synergy 21: K\"unneth Decomposition of Nested Proofs}
\label{sec:syn-nested-kunneth}

\textsc{Tags:} Tractability.

\medskip\noindent
When a proof involves \emph{nested} descent (e.g., proving a
function that branches, and within each branch calls another
function that branches), the K\"unneth formula decomposes the
composite proof:
\[
  \Hzero(\covers \times \mathcal{V})
  \cong
  \bigoplus_{p+q=0} {\Hzero}^{p}(\covers) \otimes {\Hzero}^{q}(\mathcal{V})
  = \Hzero(\covers) \otimes \Hzero(\mathcal{V})
\]

\textbf{Cubically:} the nested proof lives in a product cube
$\interval^m \times \interval^n$.  K\"unneth says this product
cube can be filled from the individual fillings of $\interval^m$
and $\interval^n$.

\textbf{Tractability gain:} instead of $m \cdot n$ fiber proofs
(one for each product fiber), we need only $m + n$ fiber proofs
(one for each factor fiber), as long as the interactions are
``tensorially simple'' (no non-trivial cross-terms in $\Hone$).


% ═══════════════════════════════════════════════════════════════════
\section{Level 6: Incremental and Caching Synergies}
\label{sec:syn-incremental}


\subsection{Synergy 22: Fiber-Granular Incremental Reverification}
\label{sec:syn-incremental-fiber}

\textsc{Tags:} Tractability.

\medskip\noindent
When code changes in one branch of a function, only the
\emph{affected fiber} needs re-verification.  The cache key is
$(\text{predicate hash}, \text{proof hash})$, and the invalidation
rule is:

\begin{enumerate}
\item Changing fiber $\varphi_k$'s code invalidates $\varphi_k$ and
  all overlaps $\varphi_k \cap \varphi_j$.
\item Changing the cover itself (adding/removing branches)
  invalidates the exhaustiveness check.
\item Unchanged fibers keep their cached verdicts.
\end{enumerate}

\textbf{Cubical reading:} changing one face of a cube invalidates
that face and all edges touching it, but not other faces.

\textbf{Cohomological reading:} this is the Mayer--Vietoris
decomposition applied to the caching structure.  The long exact
sequence:
\[
  0 \to \Hzero(X) \to \prod_i \Hzero(U_i) \to \prod_{i<j}
  \Hzero(U_i \cap U_j) \to \Hone(X) \to \cdots
\]
shows that the global section depends on local sections and
overlaps, so invalidating one local section only affects the
overlaps it participates in.


\subsection{Synergy 23: Transport-Stable Caching}
\label{sec:syn-transport-cache}

\textsc{Tags:} Tractability.

\medskip\noindent
If a proof $p$ is cached for predicate $\varphi$, and we need a
proof for predicate $\psi$ with $\psi \Rightarrow \varphi$
(Z3-verified), we can:
\begin{enumerate}
\item Look up $p$ in the cache (cache hit on $\varphi$).
\item Apply transport along $\psi \Rightarrow \varphi$ (cubical).
\item The transported proof is valid for $\psi$ without re-running Z3.
\end{enumerate}

This extends the cache to cover not just exact predicate matches
but all predicates that \emph{imply} a cached predicate ---
exponentially increasing cache hit rates for nested or refined
predicates.


% ═══════════════════════════════════════════════════════════════════
\section{Level 7: Metatheoretic Synergies}
\label{sec:syn-metatheory}


\subsection{Synergy 24: Decidability Stratification}
\label{sec:syn-decidability}

\textsc{Tags:} Tractability.

\medskip\noindent
$\mathrm{C}^4$ has a natural decidability stratification:

\begin{center}
\begin{tabular}{lll}
\toprule
\textbf{Level} & \textbf{Decidable by} & \textbf{Theory} \\
\midrule
Definitional equality & Normalization & Cubical (kernel) \\
Refinement predicates & Z3 & Cohomological (site) \\
Cover exhaustiveness & Z3 (disjunction) & Cohomological \\
Overlap emptiness & Z3 (conjunction) & Cohomological \\
Face agreement & Normalization + Z3 & Both \\
Transport validity & Z3 (implication) & Both \\
Kan filling & Algorithmic ($\Hone$ check) & Both \\
Global proof assembly & All of the above & Both \\
\bottomrule
\end{tabular}
\end{center}

Each level is decidable, but the combination is more powerful than
either theory provides alone.  The cubical kernel handles structural
equality; Z3 handles arithmetic and predicate logic; the cohomological
structure tells us \emph{which Z3 queries to issue and in what order}.


\subsection{Synergy 25: Complexity Reduction via Decomposition}
\label{sec:syn-complexity}

\textsc{Tags:} Tractability.

\medskip\noindent
Consider a function with $n$ branches, where each branch's
correctness requires a Z3 query of complexity $c$.

\begin{center}
\begin{tabular}{lcc}
\toprule
\textbf{Approach} & \textbf{Formula size} & \textbf{Z3 time} \\
\midrule
Monolithic (no decomposition) & $O(n \cdot c)$ & $O(2^{nc})$ \\
Cohomological (descent) & $n \times O(c)$ & $O(n \cdot 2^c)$ \\
+ Cubical (transport reuse) & $k \times O(c)$, $k \leq n$ &
  $O(k \cdot 2^c)$ \\
+ Overlap synthesis & $k \times O(c) + O(k^2)$ &
  $O(k \cdot 2^c + k^2)$ \\
\bottomrule
\end{tabular}
\end{center}

The cohomological decomposition reduces from $O(2^{nc})$ to
$O(n \cdot 2^c)$ (exponential in $c$ instead of $nc$).
Cubical transport reduces $n$ to $k$ (only prove dissimilar fibers,
transport to the rest).  Together, the reduction is from
exponential-in-$nc$ to polynomial-in-$k$ times exponential-in-$c$.

\begin{theorem}[Decomposition--transport theorem]
\label{thm:decomp-transport}
For a function with $n$ branches, $k$ distinct proof shapes among
them, and per-branch Z3 complexity $c$:
\[
  T_{\mathrm{C}^4}
  \;\leq\;
  k \cdot T_{\mathrm{Z3}}(c) + \binom{k}{2} \cdot T_{\mathrm{Z3}}(1)
  + (n - k) \cdot T_{\mathrm{transport}}
\]
where $T_{\mathrm{transport}} = O(1)$ (a single Z3 implication check).
\end{theorem}


\subsection{Synergy 26: Hedberg Collapse}
\label{sec:syn-hedberg}

\textsc{Tags:} Tractability.

\medskip\noindent
Hedberg's theorem: types with decidable equality are sets (all
paths between equal elements are themselves equal).

\textbf{Cohomological consequence:} for sheaves valued in sets
(our setting: all Python types have decidable equality), all
higher cohomology vanishes: $H^n = 0$ for $n \geq 1$.

\textbf{Cubical consequence:} all 2-cubes and higher are
automatically fillable (uniqueness of identity proofs).

\textbf{Synergy:} this collapse means:
\begin{enumerate}
\item We never need to check $\Htwo$ (which would require
  triple-overlap compatibility).
\item $\Hone = 0$ always, so the Kan filling condition is free.
\item The spectral sequence degenerates at $E_1$ (all higher
  differentials vanish).
\item Proof search terminates after checking fibers + pairwise
  overlaps.  No higher-dimensional complications.
\end{enumerate}

This is the reason $\mathrm{C}^4$ is \emph{tractable} for Python:
the Hedberg collapse eliminates all the ``hard'' parts of both
cubical TT and sheaf cohomology, leaving only the ``easy'' parts
(paths, covers, overlaps) that are Z3-decidable.


\subsection{Synergy 27: Propositional Truncation and Cover Independence}
\label{sec:syn-truncation}

\textsc{Tags:} Expressivity, Tractability.

\medskip\noindent
Propositional truncation $\|A\|$ turns any type into a proposition
(mere truth value).  In $\mathrm{C}^4$:

\textbf{Cubically:} truncation collapses all paths to reflexivity.
A truncated proof is ``the fact that a proof exists'' without
recording which proof.

\textbf{Cohomologically:} truncation makes the coefficient sheaf
\emph{constant} (the same value on every open).  The $\Cech$
cohomology of a constant sheaf on a paracompact space is trivial.

\textbf{Synergy:} when we only need to know that a function
satisfies a property (not how), we can truncate all fiber proofs
to $\|\mathrm{True}\|$.  This makes:
\begin{enumerate}
\item All overlaps trivially compatible (constant sheaf).
\item All transport trivial (transport of a constant is the constant).
\item The entire descent reduces to exhaustiveness checking
  ($\bigvee \varphi_i = \mathrm{True}$), which is a single Z3 query.
\end{enumerate}


% ═══════════════════════════════════════════════════════════════════
\section{Level 8: Python-Specific Synergies}
\label{sec:syn-python}


\subsection{Synergy 28: isinstance Dispatch as Fibration}
\label{sec:syn-isinstance}

\textsc{Tags:} Expressivity.

\medskip\noindent
Python's \texttt{isinstance} dispatch:
\begin{lstlisting}[language=Python]
if isinstance(x, int):    # fiber phi_int
    ...
elif isinstance(x, str):  # fiber phi_str
    ...
else:                     # fiber phi_rest
    ...
\end{lstlisting}
is simultaneously:
\begin{itemize}
\item A cubical face decomposition (three faces of a cube).
\item A refinement cover
  $\{\mathtt{isinstance}(x, \mathtt{int}),\;
  \mathtt{isinstance}(x, \mathtt{str}),\;
  \neg\mathtt{isinstance}(x, \mathtt{int}) \wedge
  \neg\mathtt{isinstance}(x, \mathtt{str})\}$.
\item A sheaf restriction: the type of \texttt{x} within
  each branch is a different fiber of the type sheaf.
\end{itemize}

\textbf{Synergy:} the compiler uses \texttt{isinstance} checks to
\emph{simultaneously} construct the refinement cover (cohomological)
and the cubical face system (cubical).  The type narrowing within
each branch ($x$ has type \texttt{int} inside the first branch)
is exactly the fiber hypothesis pushing of Synergy~6.


\subsection{Synergy 29: Dynamic Typing as Coarse-to-Fine Fibration}
\label{sec:syn-dynamic-typing}

\textsc{Tags:} Expressivity, Tractability.

\medskip\noindent
In a dynamically typed language, a variable's type is not known
statically.  In $\mathrm{C}^4$, this means the variable lives in the
\emph{total space} of the type fibration, projected to the coarsest
possible cover: $\covers = \{\mathrm{True}\}$ (a single fiber covering
all inputs).

As runtime checks and assertions narrow the type, the cover is
refined: $\{\mathrm{True}\} \to \{\texttt{isinstance}(x, T_1),\;
\neg\texttt{isinstance}(x, T_1)\} \to \cdots$

\textbf{Cubically:} the refinement adds cubical dimensions one at a
time.  Each type check adds a new face to the proof cube.

\textbf{Cohomologically:} the refinement is a sequence of cover
morphisms $\covers_0 \leftarrow \covers_1 \leftarrow \cdots$,
and Kan extension (Synergy~12) propagates proofs forward through
the refinement chain.

\textbf{Synergy:} dynamic typing is not a \emph{problem} for
$\mathrm{C}^4$ but a \emph{feature}: it provides a natural sequence
of increasingly fine covers, each of which can be verified
incrementally.


\subsection{Synergy 30: Dict-as-Record Sheafification}
\label{sec:syn-dict}

\textsc{Tags:} Expressivity.

\medskip\noindent
Python dictionaries serve as ad-hoc records.  In $\mathrm{C}^4$:

\textbf{Cohomologically:} a dictionary with keys
$\{k_1, \ldots, k_n\}$ defines a cover of the ``key space,''
and the dict is a section of the presheaf over that cover.

\textbf{Cubically:} accessing key $k_i$ is a face evaluation
($\partial_{k_i}$), and adding/removing keys is a dimensional
change.

\textbf{Synergy:} dict operations become sheaf operations:
\begin{itemize}
\item \texttt{d[k]} = restriction to the open $\{k\}$.
\item \texttt{d.update(d2)} = gluing two sections on the union
  of their key covers (requires overlap compatibility on shared keys).
\item \texttt{d.keys() == d2.keys()} = same cover.
\item Dict comprehension = a section defined locally on each key.
\end{itemize}


\subsection{Synergy 31: Exception Handling as Partial Covers}
\label{sec:syn-exceptions}

\textsc{Tags:} Expressivity.

\medskip\noindent
A \texttt{try/except} block defines a \emph{partial} cover:
\begin{itemize}
\item The \texttt{try} block is the ``normal'' fiber.
\item Each \texttt{except E} is a fiber restricted to the
  exceptional case.
\end{itemize}

\textbf{Cohomologically:} the cover $\{U_{\mathrm{normal}},
U_{E_1}, U_{E_2}, \ldots\}$ may not be exhaustive (unhandled
exceptions), and sheafification (Synergy~4) identifies the gap.

\textbf{Cubically:} the exceptional paths are \emph{alternative
faces} of the proof cube.  The proof decomposes into: normal-case
face proof + exceptional-case face proofs.

\textbf{Synergy:} the compiler can:
\begin{enumerate}
\item Check cover exhaustiveness: are all exceptions handled?
  (Z3 on the predicate ``\texttt{raises(E)}'' --- requires modeling
  exception semantics.)
\item If the cover is partial, identify the missing fiber
  (which exceptions are unhandled).
\item Transport proofs from the normal case to the exceptional case
  when the exception-handling code preserves the invariant.
\end{enumerate}


\subsection{Synergy 32: Generator/Iterator Laziness as Sheaf Sections}
\label{sec:syn-generators}

\textsc{Tags:} Expressivity.

\medskip\noindent
A Python generator yields values lazily.  In $\mathrm{C}^4$:

\textbf{Cohomologically:} the generator defines a presheaf on the
natural number site (indexed by iteration step $n$).  Each
``section'' is the value yielded at step $n$.

\textbf{Cubically:} the relationship between consecutive yields
$\mathtt{yield}[n]$ and $\mathtt{yield}[n{+}1]$ is a path (equality
or transformation proof).

\textbf{Synergy:} proving a generator correct means:
\begin{enumerate}
\item Per-step fiber proof: for each $n$, the yield value satisfies
  the specification (local section).
\item Step-to-step overlap: the transition from step $n$ to step
  $n{+}1$ is compatible (cubical path + sheaf restriction to the
  overlap $\{n\} \cap \{n{+}1\}$ in the step space).
\item Induction = descent over the natural number cover
  $\{0\}, \{1\}, \{2\}, \ldots$ with overlaps at consecutive pairs.
\end{enumerate}


% ═══════════════════════════════════════════════════════════════════
\section{Level 9: Trust and Provenance Synergies}
\label{sec:syn-trust}


\subsection{Synergy 33: Trust as Sheaf of Provenance}
\label{sec:syn-trust-sheaf}

\textsc{Tags:} Expressivity.

\medskip\noindent
The trust provenance of a proof is itself a \emph{sheaf} over the
refinement cover:
\begin{itemize}
\item Fiber $\varphi_k$ may have trust $\{\mathrm{Z3}\}$ (fully
  machine-checked).
\item Fiber $\varphi_j$ may have trust $\{\mathrm{LIBRARY}\}$
  (depends on library axiom).
\item The overlap $\varphi_k \cap \varphi_j$ inherits
  $\{\mathrm{Z3}, \mathrm{LIBRARY}\}$ (union).
\end{itemize}

\textbf{Cubically:} composing proofs with different trust levels
along a path gives the union of trust sources (trust is a
monoid under composition).

\textbf{Synergy:} the trust ``map'' over the refinement cover tells
the user \emph{exactly which parts of the input domain} rely on
which trust assumptions.  A heat map of trust over the cover:
green (Z3) on the easy fibers, yellow (library) on the hard fibers,
red (assumed) on the unverified fibers.


\subsection{Synergy 34: F*-Style Binding as Fibered Matching}
\label{sec:syn-fstar-binding}

\textsc{Tags:} Expressivity.

\medskip\noindent
F*-style binding requires that annotations match the code
structurally.  In $\mathrm{C}^4$, this is a \emph{fibered} matching:

\textbf{Cohomologically:} the annotation's branch structure must
match the code's branch structure (same refinement cover).

\textbf{Cubically:} the annotation's proof paths must connect the
same endpoints as the code's computational paths.

\textbf{Synergy:} the binding checker verifies:
\begin{enumerate}
\item Cover match: the annotation's fibers correspond 1-1 to the
  code's branches (cohomological).
\item Face match: within each fiber, the annotation's
  pre/post-conditions match the code's input/output behavior
  (cubical face evaluation).
\item Path match: the annotation's proof strategy is consistent
  with the code's control flow (cubical--cohomological: path over
  the cover).
\end{enumerate}


% ═══════════════════════════════════════════════════════════════════
\section{Summary Table}
\label{sec:syn-summary}

\begin{longtable}{clccc}
\toprule
\textbf{\#} & \textbf{Synergy} & \textbf{Expr.} & \textbf{Tract.} &
  \textbf{Auto.} \\
\midrule
\endhead
0  & Descent = HComp & $\checkmark$ & $\checkmark$ & \\
0b & Two dimension systems & $\checkmark$ & & \\
1  & Refinement types = cubical faces & $\checkmark$ & $\checkmark$ & \\
2  & Univalence in refinement lattice & $\checkmark$ & & \\
3  & Fibered universe structure & $\checkmark$ & & \\
4  & Sheafification as type completion & $\checkmark$ & & $\checkmark$ \\
5  & Transport along refinement paths & & $\checkmark$ & $\checkmark$ \\
6  & Restricted hypothesis pushing & & $\checkmark$ & \\
7  & Overlap = edge agreement & $\checkmark$ & & \\
8  & Automatic overlap synthesis & & & $\checkmark$ \\
9  & GluePath (local paths via descent) & $\checkmark$ & & \\
10 & Cover algebra = face algebra & & $\checkmark$ & \\
11 & Product covers (K\"unneth) & $\checkmark$ & $\checkmark$ & \\
12 & Cover morphisms + Kan extension & & $\checkmark$ & $\checkmark$ \\
13 & Connection-based cover simplification & & $\checkmark$ & \\
14 & Nerve topology as proof connectivity & $\checkmark$ & & \\
15 & Spectral sequence as proof search & & & $\checkmark$ \\
16 & CEGAR via cover refinement & & & $\checkmark$ \\
17 & Obstruction-guided proof repair & & & $\checkmark$ \\
18 & Transport-descent functoriality & & $\checkmark$ & $\checkmark$ \\
19 & Descent naturality & & $\checkmark$ & \\
20 & Fiber-preserving homotopies & $\checkmark$ & & \\
21 & K\"unneth for nested proofs & & $\checkmark$ & \\
22 & Fiber-granular caching & & $\checkmark$ & \\
23 & Transport-stable caching & & $\checkmark$ & \\
24 & Decidability stratification & & $\checkmark$ & \\
25 & Complexity reduction & & $\checkmark$ & \\
26 & Hedberg collapse & & $\checkmark$ & \\
27 & Propositional truncation & $\checkmark$ & $\checkmark$ & \\
28 & isinstance dispatch as fibration & $\checkmark$ & & \\
29 & Dynamic typing as coarse-to-fine & $\checkmark$ & $\checkmark$ & \\
30 & Dict-as-record sheafification & $\checkmark$ & & \\
31 & Exception handling as partial covers & $\checkmark$ & & \\
32 & Generator/iterator as sheaf sections & $\checkmark$ & & \\
33 & Trust as sheaf of provenance & $\checkmark$ & & \\
34 & F*-style binding as fibered matching & $\checkmark$ & & \\
\midrule
   & \textbf{Total} & \textbf{20} & \textbf{18} & \textbf{8} \\
\bottomrule
\caption{Complete synergy atlas.
  Expr.\ = new expressivity; Tract.\ = new tractability;
  Auto.\ = new automation.}
\label{tab:synergy-atlas}
\end{longtable}


% ═══════════════════════════════════════════════════════════════════
\section{The Synergy Theorem}
\label{sec:syn-theorem}

We conclude with the central claim of this document:

\begin{theorem}[Synergy theorem]
\label{thm:synergy}
The verification power of $\mathrm{C}^4$ is strictly greater than
the union of its cubical and cohomological components:
\[
  \mathrm{C}^4 \;\supsetneq\;
  \mathrm{CubicalTT} \;\cup\; \mathrm{SheafCoh}
\]
where $\supsetneq$ means: there exist programs whose correctness
can be stated and proved in $\mathrm{C}^4$ but not in either
cubical type theory or sheaf cohomology alone.
\end{theorem}

\begin{proof}[Proof sketch]
Consider a program $f$ with $n$ branches, where:
\begin{enumerate}
\item Branch $k$ has a correctness proof $p_k$ that requires the
  hypothesis $\varphi_k$ (the branch condition) to discharge via Z3.
\item The branches share common sub-structure, so that branches
  $i$ and $j$ have related proofs $p_i \equiv \transp_{\varphi_i
  \Rightarrow \varphi_j}(p_j)$.
\item The global proof assembles via descent + transport.
\end{enumerate}

\emph{Cubical TT alone} cannot express the decomposition into
branches (it has no site structure to define covers).  It would
need to prove the property for all inputs monolithically.

\emph{Sheaf cohomology alone} cannot express the transport
between branches (it has no path/interval structure).  It can
decompose by covers but cannot relate the local proofs to each
other except via the overlap condition.

$\mathrm{C}^4$ can: decompose by covers (cohomological), relate
branch proofs via transport (cubical), and assemble the global
proof via descent = hcomp (the foundational synergy).

The complexity reduction (Synergy~25) shows that this is not
merely a notational convenience but a genuine computational
advantage: $O(k \cdot 2^c + k^2)$ versus $O(2^{nc})$.
\end{proof}

\end{document}
